\input{../tex-inputs/latex-leseansicht-vorspann}

\section[Theodor von Sosnosky an Arthur Schnitzler, 9. 10. 1901]{L04233 Theodor von Sosnosky an Arthur Schnitzler, 9. 10. 1901}
\nopagebreak\mylabel{L04233v}
\rehead{ }\normalsize\beginnumbering\briefempfaengerindex{Schnitzler, Arthur@\textsc{Schnitzler, Arthur}!zzzSosnosky, Theodor von@\emph{von Theodor von Sosnosky}!1901-10-092@{9. 10. 1901}|(be}
\toendnotes[C]{\smallbreak\pagebreak[2]}
\correspDesc{Versand  durch Theodor von Sosnosky am 9. 10. 1901 in Alland
\newline{}Erhalt  durch Arthur Schnitzler im Zeitraum [10. 10. 1901 – 14. 10. 1901?] in Wien}\toendnotes[C]{\smallbreak}
\Standort{DLA, A:Schnitzler, HS.1985.1.4640.}
\physDesc{Brief, 1 Blatt, 4 Seiten, 1442 Zeichen
\newline{}Handschrift: blaue Tinte, lateinische Kurrent
\newline{}Schnitzler: mit rotem Buntstift eine Unterstreichung }\toendnotes[C]{\smallbreak}
\pstart
           \raggedleft{}{\pb}Alland, b. Baden\oindex{Alland@\textbf{Alland}, \emph{Verwaltungsgebiet}|pw}\pend
           
\pstart
           \raggedleft{}9./10 901\pend
           
\pstart{}Sehr geehrter Herr Doktor!\pend\vspace{0.5em}
\pstart
           Anbei erlaube ich mir, Ihnen eine \label{K_L04233-1v}\edtext{Besprechung\pwindex{Sosnosky, Theodor von 4.\,1.\,1866 Budapest – 3.\,2.\,1943 Wien@\textsc{Sosnosky, Theodor von} (4.\,1.\,1866 Budapest – 3.\,2.\,1943 Wien), \emph{Schriftsteller}!Fall Schnitzler. Eine unbefangene Betrachtung@\strich\emph{Der Fall Schnitzler. Eine unbefangene Betrachtung}|pwv}}{\lemma{\textnormal{\emph{Besprechung}}}\Cendnote{\textnormal{Theodor von Sosnosky\pwindex{Sosnosky, Theodor von 4.\,1.\,1866 Budapest – 3.\,2.\,1943 Wien@\textsc{Sosnosky, Theodor von} (4.\,1.\,1866 Budapest – 3.\,2.\,1943 Wien), \emph{Schriftsteller}|pwk}: \emph{Der Fall Schnitzler. Eine unbefangene Betrachtung}\pwindex{Sosnosky, Theodor von 4.\,1.\,1866 Budapest – 3.\,2.\,1943 Wien@\textsc{Sosnosky, Theodor von} (4.\,1.\,1866 Budapest – 3.\,2.\,1943 Wien), \emph{Schriftsteller}!Fall Schnitzler. Eine unbefangene Betrachtung@\strich\emph{Der Fall Schnitzler. Eine unbefangene Betrachtung}|pwk}. In:
                        \emph{Neue Bahnen}\pwindex{Neue Bahnen. Zeitschrift für Kunst und öffentliches Leben@\emph{Neue Bahnen. Zeitschrift für Kunst und öffentliches Leben}|pwk}, Jg. 1, Nr. 18,
                        1. 10. 1901, S. 508–514. Im
                        \emph{Österreichischen Staatsarchiv} (Signatur
                     NL-Torresani, B+C/1) werden die Korrespondenzstücke von Carl von Torresani\pwindex{Torresani-Lanzenfeld, Carl von 19.\,4.\,1846 Mailand – 16.\,4.\,1907 Nago-Torbole@\textsc{Torresani-Lanzenfeld, Carl von} (19.\,4.\,1846 Mailand – 16.\,4.\,1907 Nago-Torbole), \emph{Schriftsteller, Offizier}|pwk} an Sosnosky\pwindex{Sosnosky, Theodor von 4.\,1.\,1866 Budapest – 3.\,2.\,1943 Wien@\textsc{Sosnosky, Theodor von} (4.\,1.\,1866 Budapest – 3.\,2.\,1943 Wien), \emph{Schriftsteller}|pwk} aufbewahrt. Im Brief vom
                  8. 11. 1901 schreibt dieser: »Deine Abhandlung\pwindex{Sosnosky, Theodor von 4.\,1.\,1866 Budapest – 3.\,2.\,1943 Wien@\textsc{Sosnosky, Theodor von} (4.\,1.\,1866 Budapest – 3.\,2.\,1943 Wien), \emph{Schriftsteller}!Fall Schnitzler. Eine unbefangene Betrachtung@\strich\emph{Der Fall Schnitzler. Eine unbefangene Betrachtung}|pwv}
                     über Schnitzler habe ich mit dem grössten Interesse gelesen. Du weisst, dass
                     ich anfangs ein milderes Urtheil über die Geschichte hatte; jetzt weiss ich
                     wirklich nicht, was ich denken soll. Ich bleibe aber dabei, dass der Mann,
                     wenigstens \uline{anfänglich}, als er den ›Gustl\pwindex{Schnitzler, Arthur 15.\,5.\,1862 Wien – 21.\,10.\,1931 ebd.@\textsc{Schnitzler, Arthur} (15.\,5.\,1862 Wien – 21.\,10.\,1931 ebd.), \emph{Schriftsteller, Mediziner}!Lieutenant Gustl. Novelle@\strich\emph{Lieutenant Gustl. Novelle}|pw}‹ schrieb, nicht daran
                     dachte, den Officiersstand beleidigen zu wollen. Sowie ich Schnitzler
                     seinerzeit kennen gelernt, ist er kein Umstürzler, kein Militärfeind. Aber der
                     Erfolg von ›Freiwild\pwindex{Schnitzler, Arthur 15.\,5.\,1862 Wien – 21.\,10.\,1931 ebd.@\textsc{Schnitzler, Arthur} (15.\,5.\,1862 Wien – 21.\,10.\,1931 ebd.), \emph{Schriftsteller, Mediziner}!Freiwild. Schauspiel in 3 Akten@\strich\emph{Freiwild. Schauspiel in 3 Akten}|pw}‹, sowie der Umstand, dass er hiefür nicht zur Rechenschaft
                     gezogen wurde, mag  in seiner Auffassung die
                     Grenze zwischen dem in dieser Beziehung zulässigen und nicht Zulässigen
                     verschoben haben. Als er dann einsah, dass er sie überschritten, trat der
                     Eigensinn in Thätigkeit und erlaubte ihm nicht, einen Schritt zurückzuthun. Die
                     Ignorirung der ehrenräthlichen Vorladung dann war eine directe Tactlosigkeit
                     und Herausforderung, und kann an sich allein schon die Degradirung
                     rechtfertigen. Jedenfalls thuts mir leid, ich hatte den Mann nicht ungern
                     gehabt.« }}}\label{K_L04233-1} des vielgenannten »Lt.
                  Gustel\pwindex{Schnitzler, Arthur 15.\,5.\,1862 Wien – 21.\,10.\,1931 ebd.@\textsc{Schnitzler, Arthur} (15.\,5.\,1862 Wien – 21.\,10.\,1931 ebd.), \emph{Schriftsteller, Mediziner}!Lieutenant Gustl. Novelle@\strich\emph{Lieutenant Gustl. Novelle}|pw}« zu übersenden, den Sie so freundlich waren, mir zuzusenden.\pend
           
\pstart
           Sie werden damit zwar nicht zufrieden sein und mich – dessen bin ich gewiss – für
               einen beschränkten »Reactionär« halten, aber ich glaube anderseits, dass Sie von
               jemand, der militär-\uline{freundlich} gesinnt ist, kaum {\pb}je eine so rückhaltlose Anerkennung der künstlerischen Vorzüge des
                  Buches\pwindex{Schnitzler, Arthur 15.\,5.\,1862 Wien – 21.\,10.\,1931 ebd.@\textsc{Schnitzler, Arthur} (15.\,5.\,1862 Wien – 21.\,10.\,1931 ebd.), \emph{Schriftsteller, Mediziner}!Lieutenant Gustl. Novelle@\strich\emph{Lieutenant Gustl. Novelle}|pwv} erfahren dürften wie
               von mir. –\pend
           
\pstart
           Dass ich gegen das, was der Publication dieses Werkes\pwindex{Schnitzler, Arthur 15.\,5.\,1862 Wien – 21.\,10.\,1931 ebd.@\textsc{Schnitzler, Arthur} (15.\,5.\,1862 Wien – 21.\,10.\,1931 ebd.), \emph{Schriftsteller, Mediziner}!Lieutenant Gustl. Novelle@\strich\emph{Lieutenant Gustl. Novelle}|pwv}{ }\uline{folgte}, scharf Stellung nehme, ist eine andere
               Sache.\pend
           
\pstart
           Da ich weder im liberalen Zeitungslager, noch im gegnerischen eine Aufnahme meines
                  Artikels\pwindex{Sosnosky, Theodor von 4.\,1.\,1866 Budapest – 3.\,2.\,1943 Wien@\textsc{Sosnosky, Theodor von} (4.\,1.\,1866 Budapest – 3.\,2.\,1943 Wien), \emph{Schriftsteller}!Fall Schnitzler. Eine unbefangene Betrachtung@\strich\emph{Der Fall Schnitzler. Eine unbefangene Betrachtung}|pwv} erwarten dürfte,
               habe ich ihn einem mir sonst ganz fremden Organ\pwindex{Neue Bahnen. Zeitschrift für Kunst und öffentliches Leben@\emph{Neue Bahnen. Zeitschrift für Kunst und öffentliches Leben}|pwv} gesandt, das nicht zur Clique gehört.\pend
           
\pstart
           Ich habe acht Tage an diesem Artikel\pwindex{Sosnosky, Theodor von 4.\,1.\,1866 Budapest – 3.\,2.\,1943 Wien@\textsc{Sosnosky, Theodor von} (4.\,1.\,1866 Budapest – 3.\,2.\,1943 Wien), \emph{Schriftsteller}!Fall Schnitzler. Eine unbefangene Betrachtung@\strich\emph{Der Fall Schnitzler. Eine unbefangene Betrachtung}|pwv} gearbeitet, und i{\geminationm}er {\pb}wieder
               zugesetzt {\kaufmannsund} gestrichen, so sehr hat mich der Fall
               interessirt; und da ist es nicht unbegreiflich, dass ich ihn nicht für meine Lade
               sondern für wenigstens \uline{einige} Leser geschrieben haben
               wollte. \uline{Darum}, und nicht etwa aus materiellen
               Gründen, habe ich ihn publicirt, denn ich habe, im Voraus auf jedes Honorar
               verzichtet.\pend
           
\pstart
           Darauf gefasst, mir es mit Ihnen, Herr Doktor, für i{\geminationm}er
               verdorben zu haben, schliesse ich mit der Versicherung, dass das Urteil mir für meine
               Beziehungen zu Ihnen ganz un{\pb}massgeblich ist und dass ich in
               derselben Wertschätzung wie bisher\pend
           \pstart Hochachtungsvoll zeichne \spacefill\mbox{Th vSosnosky}\pend{}\selectlanguage{ngerman}\endnumbering\briefempfaengerindex{Schnitzler, Arthur@\textsc{Schnitzler, Arthur}!zzzSosnosky, Theodor von@\emph{von Theodor von Sosnosky}!1901-10-092@{9. 10. 1901}|)be}\mylabel{L04233h}
\begin{anhang}
\end{anhang}\newcommand{\dateiname}{L04233}\newcommand{\titel}{Theodor von Sosnosky an Arthur Schnitzler, 9. 10. 1901}\newcommand{\editorInnen}{Herausgegeben von Jahnke, SelmaMüller, Martin Anton}\input{../tex-inputs/latex-leseansicht-abspann}
